% Vietnamese translation for OI Public Library
% Source-File: IOI中国国家候选队论文/2004/李锐喆.ppt
\documentclass[11pt]{article}
\usepackage{oipl-vn}

\title{Chi tiết trong thuật toán\\{\large Ghi chú theo slide}}
\author{Li Ruizhe}
\date{IOI 2004}

\begin{document}
\maketitle

\origtitle{Details: an element that cannot be neglected}
\sourcepath{IOI 2004 / Li Ruizhe.ppt}

\section{Vai trò của chi tiết}

Tệp PPT này chứa nhiều công thức dưới dạng đối tượng nhúng, nên phần văn bản
trích được không đầy đủ. Nội dung khớp với bài viết PDF: trong thi thuật
toán, chi tiết cài đặt không phải phần phụ; nó có thể quyết định cả độ đúng
lẫn độ phức tạp thực tế.

\section{Hai loại chi tiết}

Loại thứ nhất gần như không đổi bậc độ phức tạp nhưng làm chương trình gọn,
dễ chứng minh và ít lỗi hơn. Ví dụ là cách duy trì danh sách, cách đánh dấu,
cách chọn mảng phụ hoặc thứ tự cập nhật.

Loại thứ hai tác động trực tiếp đến bậc thời gian. Khi một thao tác nhỏ nằm
trong vòng lặp sâu, việc dùng cấu trúc dữ liệu phù hợp, nén quan hệ hoặc lưu
kết quả trung gian có thể biến một lời giải không chạy kịp thành lời giải
đạt yêu cầu. Bản dịch PDF phân tích các ví dụ về quan hệ nhóm, đường sắt và
những thao tác cần được tổ chức cẩn thận.

\section{Ghi nhớ}

Thông điệp của slide là: sau khi có ý tưởng thuật toán, phải kiểm tra từng
điểm nối giữa mô hình và chương trình. Biên mảng, thứ tự cập nhật, điều kiện
dừng, cách so sánh và cách khôi phục trạng thái đều là một phần của thuật
toán, không phải việc trang trí sau cùng.

\end{document}
